% Phase V + VI empirical section.
% Self-contained LaTeX fragment for insertion into
% paper/neurips_selective_temporal_credit_assignment_positioned.tex.
% References expect: \texttt{graphicx}, \texttt{booktabs}, \texttt{amsmath},
% \texttt{amssymb}, \texttt{hyperref}, and \texttt{cleveref}.

\section{Experiments}
\label{sec:experiments}

This section presents the Phase V + VI empirical program. In
\Cref{ssec:phaseVI_design} we state the mechanism-first experimental
principles. In \Cref{ssec:phaseVI_families} we describe five constructed
task families and the fixed promotion gate used to select shortlists.
\Cref{ssec:phaseVI_results} lays out the main-text evidence:
(i) the safe operator changes the optimal policy at near-tied contest
states (Family A); (ii) it produces a \emph{Pareto improvement} on the
dual-objective surface under stochastic transitions; (iii) clipping
preserves the contraction certificate on visited transitions where the
raw weighted-LSE operator violates it (Family C). Phase I--IV appear in
\Cref{app:phase_I_IV_sanity} as sanity material only.

\subsection{Experimental principles}
\label{ssec:phaseVI_design}

Phase V is a theorem-first rebuild: we design tasks whose classical
decision boundary is \emph{intentionally near-indifferent}, so the
nonlinear operator has room to flip the preferred action. Promotion to
the main-text shortlist requires the following hard-coded thresholds,
which are never relaxed:
(a) $\mathrm{policy\_disagreement} \geq 0.05$ under the reference
occupancy $d_\mathrm{ref} = \tfrac12 d^{\pi^*_\mathrm{cl}} + \tfrac12 d^{\pi^*_\mathrm{safe}}$
\emph{or} start-state greedy action differs;
(b) $\mathrm{mass\_delta\_d} \geq 0.10$;
(c) $|vgn| := |V^*_\mathrm{safe}(s_0) - V^*_\mathrm{cl}(s_0)|/R_\mathrm{scale} \geq 0.005$;
(d) $\mathrm{contest\_gap\_norm} \leq 0.01$ (strict near-tie band); and
(e) no certification violation on visited mass.

For clipping activity we admit two regimes: $\texttt{binding\_clip}$
($0.05 \leq \mathrm{clip\_fraction} \leq 0.80$) and
$\texttt{safe\_active\_no\_distortion}$ ($\mathrm{clip\_fraction}=0$ with
the raw weighted-LSE operator already inside the certified box on
visited transitions). The latter captures a mechanism-relevant regime
the original search criterion discarded (see \S\ref{ssec:phaseVI_results}).
Family~C (the safety family) retains a binding-clip requirement.

\subsection{Task families}
\label{ssec:phaseVI_families}

Five constructed families cover a principled grid: \emph{Family~A}
(delayed jackpot vs smooth stream, concentration contrast);
\emph{Family~B} (rare catastrophe vs safe branch, alternative-branch
propagation); \emph{Family~C} (raw-stress, safety);
\emph{Family~D} (preventive intervention, propagation-depth warning);
\emph{Family~E} (regime-shift with concentration contrast). Each family
exposes a tie parameter $\lambda$ (a classical-near-indifference knob)
and a geometry parameter $\psi$. The search driver enumerates
$\le 5{,}000$ candidates per iteration, solves $\lambda_\mathrm{tie}(\psi)$
via bisection around a closed-form warm start, sweeps $\varepsilon$-bands
around $\lambda_\mathrm{tie}$, and evaluates every $\S 6$ metric
(\Cref{tab:phaseVI_summary}).

\paragraph{Key structural finding.} Families B and D sweep over 1,800
candidates combined and produce \emph{zero} promotable tasks; Family E
sweeps 540 and also promotes zero. The diagnosis holds uniformly: in
these \emph{sequential-chain} constructions (action at $t=0$ gates a
downstream branch structure), value translation accumulates too slowly
along the chain to clear the 5e-3 $|vgn|$ threshold. Family A's
\emph{side-by-side contest} construction (both branches propagate
simultaneously from the same contest state) clears it readily. The
design insight is mechanism-theoretic: the safe weighted-LSE operator's
value translation is strong when candidate branches are evaluated
\emph{simultaneously} from the same backward-sweep step, and weak when
the decision gates a chain whose reward geometry is realized in only
one branch. B, D, E are therefore appendix-only design diagnostics; the
main text is carried by A and C.

\subsection{Main-text evidence}
\label{ssec:phaseVI_results}

\paragraph{Mechanism: the safe operator changes the optimal policy.}
On the 5 deterministic Family~A shortlisted tasks (all promotion mode
$\texttt{safe\_active\_no\_distortion}$), the safe operator produces a
start-state policy flip and $|vgn|$ up to $7.82\times 10^{-3}$. Exact
greedy recovery under limited backups shows safe reaching full policy
correctness at $k = T-1$ while classical reaches it at $k = T$
(\Cref{fig:propagation_curve,fig:stagewise_heatmap,fig:greedy_recovery}).
This one-step-earlier propagation is consistent with the theorem's
prediction that on aligned regions the safe operator propagates a
delayed signal backward differently from classical discounting.

\paragraph{Dual-objective Pareto (novel).}
On the stochastic Family~A shortlist (6 tasks, $p_\mathrm{transit}=0.95$,
$\pi^*_\mathrm{safe}\neq \pi^*_\mathrm{cl}$ at $s_0$ on every task), we
compute exact policy evaluation under both objectives:
\[
\underbrace{V^{\pi^*_\mathrm{cl}}_\mathrm{cl}(s_0) - V^{\pi^*_\mathrm{safe}}_\mathrm{cl}(s_0)}_{\text{classical-eval regret}} \approx 0 \quad\text{and}\quad
\underbrace{V^{\pi^*_\mathrm{safe}}_\mathrm{safe}(s_0) - V^{\pi^*_\mathrm{cl}}_\mathrm{safe}(s_0)}_{\text{safe-eval advantage}} > 0.
\]
Under classical evaluation, the safe policy matches classical return
(no regret) to within ULP ($\le\!4.4\times 10^{-16}$ across all 6
tasks). Under the safe (nonlinear weighted-LSE) evaluator, the safe
policy strictly dominates, with gap scaling linearly from $0.012$ at
$R=2$ to $0.499$ at $R=100$ (\Cref{fig:dual_objective_eval}). The safe
policy is therefore \emph{Pareto-better on the dual-objective surface}.

This reframes the RL translation question: any return-AUC comparison
under a single (classical) metric measures zero by the classical-tie
construction; the mechanism lives in the \emph{objective}, not in a
sample-noise advantage.

\paragraph{Risk profile (VI-G): the safe operator is a signed tie-breaker, not a universal risk-reducer.}
Ten-thousand Monte-Carlo rollouts per policy per task produce return
distributions from which we compute variance and conditional
value-at-risk $\mathrm{CVaR}_\alpha$ at $\alpha \in \{0.05, 0.10, 0.25\}$.
On Family~A's optimistic sign branch ($\mathrm{sign\_family}=+1$):
(i) $\mathbb{E}[G | \pi^*_\mathrm{safe}] \approx \mathbb{E}[G | \pi^*_\mathrm{cl}]$ (matches the dual-eval result);
(ii) $\mathrm{Var}[G | \pi^*_\mathrm{safe}] \gg \mathrm{Var}[G | \pi^*_\mathrm{cl}]$ (safe picks the delayed-jackpot branch whose return distribution is bimodal);
(iii) $\mathrm{CVaR}_{0.05}[G | \pi^*_\mathrm{safe}] < \mathrm{CVaR}_{0.05}[G | \pi^*_\mathrm{cl}]$ (safe has worse worst-case).

This is theoretically consistent: with optimistic sign, the safe
operator amplifies gains rather than penalizing losses, producing
risk-seeking credit assignment. The pessimistic sign (Family~B/D
failure mode) would symmetrically produce a CVaR-improving safe
policy, had B or D promoted tasks. The correct interpretation is that
the safe weighted-LSE operator is a \emph{signed, theorem-linked
tie-breaker between statistically-equivalent policies}, with the sign
determining whether the correction favors gain amplification or loss
penalization. It is not a universal risk-reducer.

\paragraph{Safety: clipping preserves the contraction cert.}
On the stronger Family~C shortlist (84 tasks, $\beta_\mathrm{raw}$ up to
128$\times$ the certified cap), the safe-clipped operator's local
continuation derivative stays within the contraction certificate,
$d_\mathrm{safe} \leq \kappa_t \in [0.95, 0.96]$ on every visited
transition. The raw unclipped operator \emph{violates} the certificate,
with $d_\mathrm{raw}$ up to $1.416$ on visited mass (a $48\%$
amplification margin over the classical discount). Classifying each
task's operator sweep by the 5-label taxonomy $\{$\texttt{stable\_linear},
\texttt{stable\_nonlinear}, \texttt{expansive\_bounded},
\texttt{expansive\_unbounded}, \texttt{nan\_guarded}$\}$, classical is
\texttt{stable\_linear} on all $84$ tasks, safe is \texttt{stable\_nonlinear}
on all $84$, and raw is \texttt{expansive\_bounded} on all $84$.
\Cref{fig:safe_vs_raw_stability} visualizes the per-stage $V$
trajectories and the $d_\mathrm{raw}$ histogram under $d_\mathrm{ref}$.

Finite-horizon backward sweeps truncate amplification, so all three
operators produce bounded $V$ at horizon $T$. To expose the contraction
difference directly we iterate the backward-sweep operator $F\colon
V_T \mapsto V_0$ for 50 outer steps with perturbed $V_T$ boundary
(\Cref{fig:perturbation_growth}): even raw operators that violate
$d \leq \kappa_t$ per-transition retain a bounded fixed point of $F$.
The certificate bound is therefore \emph{sufficient} for contraction
but not \emph{necessary} for finite-horizon $V$-boundedness; what
clipping prevents is loss of the theoretical contraction guarantee
(required for infinite-horizon or adversarial extensions), not immediate
functional divergence.

\paragraph{What the RL experiments did and did not show.}
A full 5-arm Stage-1 paired RL pilot (classical, safe-zero,
safe-nonlinear, tuned fixed-$\gamma$, TD($\lambda$), 2 tasks $\times$ 5
seeds) produces identical AUC across every arm under classical-return
evaluation ($\mathrm{CI}_{95\%} = [0, 0]$ exactly). A follow-up
DP-$Q$-initialized fine-tune preserves the DP-optimal start-state argmax
on 100\% of runs but still measures AUC $\equiv 0$. The dual-evaluation
result above explains this structurally: at the classical tie, both
branches achieve equal classical return by construction, so any classical-
return RL evaluator measures zero regardless of sample budget. The
paper's RL claim is therefore intentionally scoped: the mechanism is
\emph{planner}-level, and our RL null confirms the theory's scope
prediction from \Cref{sec:near_zero_beta_regime} rather than contradicting
the mechanism.

\subsection{Consistency audit}
\label{ssec:audit}

A Phase V consistency audit (\texttt{scripts/audit/run\_consistency\_audit.py})
recomputes every Phase I--IV table and figure from raw logs and
cross-checks against paper quantities. The first-pass audit flagged 82
DP cert-bound BLOCKER rows on Phase III \texttt{SafePE}/\texttt{SafeVI}
runs; the remediation pass (occupancy-weighted cert check,
\texttt{scripts/audit/occupancy\_cert\_check.py}) demonstrated that
these are benign artifacts of the P-weighted effective-discount
aggregation on absorbing cells, and downgraded all 82 to MINOR. One
paper-text discrepancy remains: a stale $79\%$ \texttt{clip\_fraction}
number for \texttt{grid\_hazard} (recomputed: $96.25\%$); we correct
this figure in the appendix and flag it as a consistency-pipeline
outcome.

\begin{table*}[t]
\centering
\small
\begin{table}[t]
\centering
\small
\caption{Headline empirical claims. All numbers are from exact
finite-horizon planning / exact policy evaluation on tabular MDPs;
no RL learning curves are included in the main text. Abbreviations:
$|vgn|=|V^*_{\mathrm{safe}}(s_0)-V^*_{\mathrm{cl}}(s_0)|/R_{\mathrm{scale}}$;
\emph{PD}, policy disagreement under $d_{\mathrm{ref}}$;
\emph{flip rate}, fraction of shortlisted tasks on which the
safe-optimal and classical-optimal greedy actions differ at $s_0$.}
\label{tab:phaseVI_summary}
\begin{tabular}{l l r l}
\toprule
Claim & Evidence & Tasks & Headline numbers \\
\midrule
Objective-level
  & deterministic Family~A & $5$
  & max $|vgn|=7.82\!\times\!10^{-3}$,\,max PD $=0.25$,\,flip rate $=1.0$ \\
policy improvement
  & (near-tie, \texttt{safe\_active\_no\_distortion}) & & \\
\addlinespace[3pt]
Dual-objective
  & stochastic Family~A & $6$
  & classical-eval regret $\le 4.4\!\times\!10^{-16}$; \\
Pareto improvement
  & (near-tie, $p_{\mathrm{transit}}=0.95$) & &
  safe-eval advantage $\in[1.2\!\times\!10^{-2},\,4.99\!\times\!10^{-1}]$, \\
  & & & mean $\approx 0.156$, flip rate $=1.0$ \\
\addlinespace[3pt]
Certified stability
  & Family~C raw-stress & $84$
  & safe $d_t\le\kappa_t\in[0.95,0.96]$;\,raw $d_t$ up to $1.416$; \\
  & (\texttt{binding\_clip}) & &
  safe/classical $V[0,s_0]\!\approx\!4.92/4.86$; raw $V[0,s_0]\!\approx\!16.9$ \\
\bottomrule
\end{tabular}
\end{table}

\end{table*}

%% Figures (main-text; files under figures/main/):
%% - fig_propagation_curve.pdf         (WP3 Family A greedy recovery)
%% - fig_stagewise_error_heatmap.pdf   (WP3 stagewise error)
%% - fig_greedy_recovery.pdf           (WP3 greedy correctness)
%% - fig_safe_vs_raw_stability.pdf     (WP4 per-stage V + derivative histogram)
%% - fig_perturbation_growth.pdf       (VI-C iterated-F stability)
%% - fig_dual_objective_eval.pdf       (VI-F dual-objective Pareto)

%% Appendix (C + sanity):
%% - fig_phase123_sanity.pdf    (Phase I--IV benchmark matrix)
%% - fig_appendix_ablations.pdf (Family B/D/E design diagnostics)
